\section{Recommandations et mise en œuvre}\label{sec-reco}

\begin{objectif}
Traduire les résultats en règles de contrôle, préciser les données nécessaires et la procédure mensuelle.
\end{objectif}

\subsection{Ce que nous recommandons}
\begin{enumerate}
\item \textbf{Contrôler le forward, pas la répartition.} Pour chaque portefeuille, calculer la ligne forward
(forwards Totem, dividendes FO) et la ligne non-forward $\sprod{\eta}{\Delta\ell}$. Pour les portefeuilles linéaires en
forward, la ligne forward \emph{est} le contrôle, et elle ne dépend pas des contributeurs dividendes
(\cref{thm-control,thm-exact}). Les dividendes Totem ne sont nécessaires que là où $\eta\ne0$ : swaps et futures de
dividendes, prêt-emprunt de titres, financement des swaps de rendement total, exercice américain, exotiques.
\item \textbf{Ne pas additionner sans compensation} des lignes dividendes et repo importantes et de signes opposés :
c'est un artefact de répartition. En valorisation prudente, prendre le forward comme paramètre là où $\eta\approx0$.
\item \textbf{Juger la courbe en portage, avec une bande de bruit.} Remplacer les contrôles visuels et les seuils en
points de base de repo par l'écart de portage $\delta$ et la covariance $\Gamma$ calculée à partir des dispersions et
des nombres de contributeurs Totem (\cref{prop-noise}).
\item \textbf{Déployer le détecteur structurel} avec $\alpha=\pct{2}$, et publier pour chaque alerte l'explication en
unités naturelles, la courbe réparée et, à part, le test d'erreur de marquage FO.
\item \textbf{Garder les B-splines, les mélanges et les réseaux de neurones en aval} : sur les sorties du détecteur et
ses résidus, et pour regrouper la calibration par région et catégorie.
\item \textbf{Calibrer et valider sur l'historique réel} : $\kappa$ et l'échelle de niveau par région et catégorie,
chaque trimestre, puis une validation sur douze à vingt-quatre mois d'instantanés avec des incidents étiquetés.
\end{enumerate}

\subsection{Données nécessaires}
\begin{center}\small
\begin{tabular}{@{}>{\raggedright\arraybackslash}p{2.6cm}>{\raggedright\arraybackslash}p{9.3cm}
>{\raggedright\arraybackslash}p{2.6cm}@{}}
\toprule
Donnée & Champs & Taille\\
\midrule
Options Totem & maturités ; strikes ; prix de consensus des calls et des puts ; nombre de contributeurs ; dispersion des
forwards des contributeurs (ou une estimation) ; spot de référence & $N$ maturités\\
Dividendes Totem & bornes des périodes ; niveau de consensus ; dispersion ; nombre de contributeurs & $m$ périodes\\
Marché FO & spot à l'heure de l'instantané Totem ; courbe des taux ; repo aux maturités Totem ; prévision de
dividendes (dates et montants) & $N$ ; liste\\
Risque FO & sensibilités repo par pilier, projetées sur les maturités Totem ; sensibilités dividendes par période
Totem & $N+m$ par portefeuille\\
Historique & douze à vingt-quatre instantanés mensuels et incidents connus & calibration, validation\\
\bottomrule
\end{tabular}
\end{center}
Si Totem ne publie pas la dispersion des forwards des contributeurs, on peut l'estimer à partir de la dispersion des
calls et des puts très dans la monnaie, ou la calibrer sur l'historique : le détecteur a besoin de la bonne forme par
maturité, et $\kappa$ corrige un biais commun.

\subsection{Procédure mensuelle}
\begin{enumerate}
\item Charger les instantanés Totem et FO à la même heure de valorisation ; vérifier le spot de référence.
\item Calculer les forwards par la parité ; examiner le facteur d'actualisation implicite, les résidus de parité et
les rapports de dispersion.
\item Allouer les niveaux de dividendes Totem sur le profil de détachement FO ; calculer le repo implicite avec les
dividendes Totem ($q^{\tot}$) et avec les dividendes FO ($\hat q$).
\item Construire $\Gamma$ et lancer le détecteur avec le $\kappa$ et l'échelle de niveau de la région et de la catégorie.
\item Calculer les lignes de contrôle par portefeuille (forward, dividendes, repo, non-forward) et la matérialité de la
répartition.
\item Examiner les alertes avec leurs explications, ainsi que les erreurs de marquage FO possibles ; utiliser la courbe
réparée quand une anomalie Totem est confirmée.
\item Consigner la décision prise pour chaque alerte : ces étiquettes servent à la calibration suivante et à la
validation.
\end{enumerate}

\subsection{Gouvernance et prochaines étapes}
$\kappa$ et l'échelle de niveau sont recalibrés chaque trimestre par région et catégorie sur l'historique sain, avec un
budget d'alertes de \vFPR{} des sous-jacents par mois. On suit le taux d'alerte réalisé et la part des explications
acceptées par les contrôleurs. Là où elles existent, des références de marché complètent le détecteur : contrats à
terme cotés sur indices et sur dividendes, contrats à terme de rendement total, commissions de prêt-emprunt de titres
pour le court terme. L'étape suivante est de faire tourner la chaîne sur douze à vingt-quatre mois d'instantanés Totem
et FO réels, de calibrer les réglages par région et catégorie, et de mesurer la qualité des explications sur des
incidents étiquetés.

\section*{Conclusion}
\addcontentsline{toc}{section}{Conclusion}
Les courbes de repo implicite sont irrégulières parce qu'elles sont le résidu de deux moyennes de consensus calculées
sur des panels différents, et une partie de ce résidu ne peut être retirée par aucune méthode qui ne voit que le
consensus (\cref{thm-nonid}). Deux conséquences pratiques en découlent. Pour le contrôle, un portefeuille linéaire en
forward n'a pas besoin de la répartition entre repo et dividendes : le contrôle repo est le contrôle forward moins le
contrôle dividendes, exactement en réévaluation complète, et l'exposition non-forward $\eta$ dit, portefeuille par
portefeuille, quand la répartition compte. Pour la détection, les mécanismes qui produisent les vraies anomalies ont
des formes connues dans l'écart de portage ; des filtres adaptés, avec un bruit calculé à partir des propres
statistiques de Totem, les détectent à un taux de fausses alertes fixé, les expliquent en unités naturelles et les
distinguent des erreurs de marquage FO.
